% !TeX root = ../main.tex

\section{Implementazione}
L'implementazione del database è stata organizzata separando il caricamento dei dati, la logica procedurale lato database e gli script di test delle principali operazioni. In particolare:
\begin{itemize}
  \item lo script \texttt{04\_data.sql} gestisce il popolamento della base di dati a partire da file CSV;
  \item lo script \texttt{03\_logic.sql} definisce funzioni e trigger \texttt{PL/pgSQL} per i controlli dinamici e gli aggiornamenti automatici;
  \item lo script \texttt{05\_queries.sql} raccoglie una serie di interrogazioni e operazioni di prova utilizzate per verificare il comportamento del sistema.
\end{itemize}

\subsection{Generazione dei dati di test con Faker}
Per disporre di una base di dati iniziale sufficientemente ampia e realistica è stato sviluppato uno script Python dedicato, \texttt{generate\_data.py}, basato sulla libreria \texttt{Faker} con localizzazione italiana. Lo script produce automaticamente file CSV contenenti dati plausibili per il dominio applicativo, come nomi, indirizzi, codici fiscali, date e importi.

La generazione è parametrica: il numero di banche, filiali, persone, clienti, impiegati, conti, carte Bancomat, interfacce e operazioni può essere modificato tramite costanti definite all'inizio dello script. Questo approccio ha permesso di creare rapidamente insiemi di dati coerenti con la struttura del database e adeguati al collaudo delle operazioni implementate.

Durante la generazione si è posta attenzione soprattutto alla coerenza referenziale tra i dati prodotti. In particolare, i conti vengono associati a filiali esistenti, le carte Bancomat vengono collegate a clienti e conti già generati, la storia dei servizi degli impiegati viene costruita in modo sequenziale e le operazioni vengono prodotte in modo compatibile con il tipo di interfaccia utilizzata.

\subsection{Popolamento}
Il popolamento del database è stato realizzato tramite lo script \texttt{04\_data.sql}, che carica i dati da file CSV mediante il comando \texttt{\textbackslash copy}. Poiché alcune tabelle utilizzano colonne definite con \texttt{GENERATED ALWAYS AS IDENTITY}, prima del caricamento tali colonne vengono temporaneamente impostate come \texttt{GENERATED BY DEFAULT}, così da consentire l'inserimento esplicito degli identificativi presenti nei file.

\begin{lstlisting}[caption={Gestione temporanea delle colonne IDENTITY}]
ALTER TABLE banca ALTER COLUMN codiceBanca SET GENERATED BY DEFAULT;
ALTER TABLE filiale ALTER COLUMN codiceFiliale SET GENERATED BY DEFAULT;
ALTER TABLE impiegato ALTER COLUMN idImpiegato SET GENERATED BY DEFAULT;
ALTER TABLE conto ALTER COLUMN numeroConto SET GENERATED BY DEFAULT;
ALTER TABLE interfaccia ALTER COLUMN idInterfaccia SET GENERATED BY DEFAULT;
ALTER TABLE operazione ALTER COLUMN idOperazione SET GENERATED BY DEFAULT;
\end{lstlisting}

Durante l'importazione bulk vengono inoltre disabilitati i trigger sulle tabelle \texttt{presta\_servizio\_in} e \texttt{operazione}, in modo da evitare controlli e side-effect sui dati storici caricati dai file.

\begin{lstlisting}[caption={Disabilitazione temporanea dei trigger durante il caricamento}]
ALTER TABLE presta_servizio_in DISABLE TRIGGER ALL;
ALTER TABLE operazione DISABLE TRIGGER ALL;

\copy banca FROM '${DATA_DIR}/banca.csv' WITH (FORMAT csv, NULL '')
\copy filiale FROM '${DATA_DIR}/filiale.csv' WITH (FORMAT csv, NULL '')
\copy persona FROM '${DATA_DIR}/persona.csv' WITH (FORMAT csv, NULL '')
\copy cliente FROM '${DATA_DIR}/cliente.csv' WITH (FORMAT csv, NULL '')
\copy impiegato FROM '${DATA_DIR}/impiegato.csv' WITH (FORMAT csv, NULL '')
\copy conto FROM '${DATA_DIR}/conto.csv' WITH (FORMAT csv, NULL '')
\copy intestato_a FROM '${DATA_DIR}/intestato_a.csv' WITH (FORMAT csv, NULL '')
\copy presta_servizio_in FROM '${DATA_DIR}/presta_servizio_in.csv' WITH (FORMAT csv, NULL '')
\copy cartaBancomat FROM '${DATA_DIR}/cartaBancomat.csv' WITH (FORMAT csv, NULL '')
\copy interfaccia FROM '${DATA_DIR}/interfaccia.csv' WITH (FORMAT csv, NULL '')
\copy operazione FROM '${DATA_DIR}/operazione.csv' WITH (FORMAT csv, NULL '')

ALTER TABLE presta_servizio_in ENABLE TRIGGER ALL;
ALTER TABLE operazione ENABLE TRIGGER ALL;
\end{lstlisting}

Al termine del caricamento vengono ripristinate le colonne \texttt{IDENTITY} nella modalità originaria \texttt{GENERATED ALWAYS} e vengono sincronizzate le sequenze interne con il valore massimo effettivamente presente nelle tabelle, così da evitare collisioni nei successivi inserimenti.

\begin{lstlisting}[caption={Ripristino delle sequenze dopo il caricamento}]
ALTER TABLE banca ALTER COLUMN codiceBanca SET GENERATED ALWAYS;
ALTER TABLE filiale ALTER COLUMN codiceFiliale SET GENERATED ALWAYS;
ALTER TABLE impiegato ALTER COLUMN idImpiegato SET GENERATED ALWAYS;
ALTER TABLE conto ALTER COLUMN numeroConto SET GENERATED ALWAYS;
ALTER TABLE interfaccia ALTER COLUMN idInterfaccia SET GENERATED ALWAYS;
ALTER TABLE operazione ALTER COLUMN idOperazione SET GENERATED ALWAYS;

SELECT setval(pg_get_serial_sequence('banca', 'codicebanca'),
    (SELECT MAX(codiceBanca) FROM banca));
SELECT setval(pg_get_serial_sequence('filiale', 'codicefiliale'),
    (SELECT MAX(codiceFiliale) FROM filiale));
SELECT setval(pg_get_serial_sequence('impiegato', 'idimpiegato'),
    (SELECT MAX(idImpiegato) FROM impiegato));
SELECT setval(pg_get_serial_sequence('conto', 'numeroconto'),
    (SELECT MAX(numeroConto) FROM conto));
SELECT setval(pg_get_serial_sequence('interfaccia', 'idinterfaccia'),
    (SELECT MAX(idInterfaccia) FROM interfaccia));
SELECT setval(pg_get_serial_sequence('operazione', 'idoperazione'),
    (SELECT MAX(idOperazione) FROM operazione));
\end{lstlisting}

\subsection{Trigger}
La parte procedurale del progetto è stata implementata mediante funzioni e trigger \texttt{PL/pgSQL}, definiti nello script \texttt{03\_logic.sql}. Tali componenti completano i vincoli statici dello schema fisico con controlli dinamici e aggiornamenti automatici.

Le funzionalità introdotte sono le seguenti:
\begin{itemize}
  \item \textbf{check\_correttezza\_carta}: verifica, prima dell'inserimento di una operazione, che la carta Bancomat esista e sia associata al conto indicato;
  \item \textbf{check\_trasferimento\_insert}: controlla che un impiegato non venga assegnato a una filiale presso cui ha già prestato servizio e che sia trascorso almeno un mese dall'ultima assegnazione;
  \item \textbf{calcola\_commissione\_prelievo}: funzione di supporto che determina la commissione da applicare a un prelievo in base alla banca del conto e a quella dell'interfaccia utilizzata;
  \item \textbf{check\_prelievo}: blocca i prelievi non validi, verificando saldo del conto, disponibilità dell'interfaccia e limiti giornalieri e mensili della carta;
  \item \textbf{update\_post\_operazione}: aggiorna automaticamente, dopo l'inserimento di una operazione, il numero di operazioni dell'interfaccia, il saldo del conto, la disponibilità dell'interfaccia e le soglie di spesa della carta;
  \item \textbf{check\_conto\_estinto}: impedisce l'inserimento di operazioni riferite a conti già estinti;
  \item \textbf{gestione\_conto\_estinto}: gestisce l'estinzione di un conto corrente, azzerandone il saldo (impostato a NULL) e rimuovendo le relative carte Bancomat e operazioni collegate per ragioni di sicurezza;
  \item \textbf{check\_operazione\_cassa}: assicura l'integrità operativa delle casse fisiche, impedendo che vi vengano registrate operazioni diverse da prelievi e depositi.
\end{itemize}

I trigger associati a queste funzioni sono definiti come segue:

\begin{lstlisting}[caption={Definizione dei trigger principali}]
create trigger trg_correttezza_carta before insert
    on consorzio.operazione
    for each row execute function consorzio.check_correttezza_carta();

create trigger trg_check_trasferimento_insert before insert
    on consorzio.presta_servizio_in
    for each row execute function consorzio.check_trasferimento_insert();

create trigger trg_blocco_prelievo before insert
    on consorzio.operazione
    for each row execute function consorzio.check_prelievo();

create trigger trg_update_post_op after insert
    on consorzio.operazione
    for each row execute function consorzio.update_post_operazione();

create trigger trg_check_conto_estinto before insert
    on consorzio.operazione
    for each row execute function consorzio.check_conto_estinto();

create trigger trg_gestione_conto_estinto before insert or update
    on consorzio.conto
    for each row execute function consorzio.gestione_conto_estinto();

create trigger trg_check_operazione_cassa before insert
    on consorzio.operazione 
    for each row execute function consorzio.check_operazione_cassa();
\end{lstlisting}

Nel complesso, questa logica consente di mantenere centralizzati nel database sia i controlli di validazione sia gli aggiornamenti conseguenti alle operazioni effettuate.

\subsection{Operazioni e Query}
Per verificare il corretto funzionamento del sistema è stato predisposto lo script \texttt{05\_queries.sql}, che raccoglie una serie di operazioni e interrogazioni rappresentative dei casi d'uso principali del progetto.

Le prove implementate sono le seguenti:
\begin{enumerate}
  \item estrazione delle informazioni relative a una specifica interfaccia;
  \item inserimento di una operazione di prelievo su un conto esistente, con attivazione dei controlli e degli aggiornamenti automatici lato database;
  \item inserimento di un nuovo impiegato e contestuale assegnazione iniziale a una filiale mediante una sequenza di \texttt{CTE};
  \item trasferimento di un impiegato verso una nuova filiale, seguito dall'aggiornamento della precedente assegnazione;
  \item licenziamento di un impiegato mediante cancellazione del relativo record;
  \item individuazione del cliente con ammontare complessivo più elevato sui conti non estinti a lui intestati;
  \item individuazione delle filiali per cui tutte le operazioni registrate sulle rispettive interfacce hanno ammontare strettamente maggiore di 20;
  \item ricerca degli impiegati che sono anche clienti e che risultano titolari di almeno un conto aperto da almeno 13 anni;
  \item ricerca delle coppie di clienti che risultano intestatarie di conti presso esattamente lo stesso insieme di filiali, con almeno due filiali distinte.
\end{enumerate}

A titolo esemplificativo, l'inserimento di un nuovo impiegato viene realizzato concatenando più operazioni in un'unica istruzione, così da creare prima la persona, poi il relativo impiegato e infine l'assegnazione iniziale in filiale.

\begin{lstlisting}[caption={Inserimento di un nuovo impiegato e assegnazione iniziale}]
with new_user as (
    insert into Persona (codiceFiscale, nome, indirizzo) values
    ('BNCMRR92L10H501T', 'Bianca Rossi', 'Via delle Scienze 208, Udine')
    on conflict (codiceFiscale) do update
        set codiceFiscale = EXCLUDED.codiceFiscale
    returning codiceFiscale
),
new_employee as (
    insert into Impiegato (dataAssunzione, codiceFiscale_i)
    select CURRENT_DATE, codiceFiscale
    from new_user
    returning idImpiegato
)
insert into presta_servizio_in (idImpiegato_psi, codiceFiliale_psi, dataInizio, dataFine)
select idImpiegato, 1, CURRENT_DATE, null
from new_employee;
\end{lstlisting}

Un secondo esempio significativo è dato dalla query che individua il cliente con il totale più alto sui propri conti attivi.

\begin{lstlisting}[caption={Cliente con ammontare complessivo più elevato}]
select
    p.codiceFiscale,
    p.nome,
    p.indirizzo,
    sum(co.ammontare)
from persona as p
join cliente as c on p.codicefiscale = c.codiceFiscale_c
join intestato_a as ia on c.codiceFiscale_c = ia.codiceFiscale_int
join conto as co on ia.numeroConto_int = co.numeroConto
where co.dataEstinzione is null
group by p.codiceFiscale, p.nome, p.indirizzo
order by sum(co.ammontare) desc
limit 1;
\end{lstlisting}

Queste interrogazioni sono state utilizzate sia per verificare la correttezza della base di dati sia per collaudare il comportamento congiunto di vincoli, trigger e aggiornamenti automatici.
